%% kapa: exam=1 round=34 year=2023 period=2 form=A subject=appraisal_law subject_ko=감정평가관계법규 seq=14 number=14
%% source: Q-Net 감정평가사 기출문제 (cst00309) artlSeq=5211789 file=raw/1008/5211789/2023년도 제34회 감정평가사 1차시험 2교시.pdf
%% conversion: kordoc 4.24.51 -> tools/gen1.py (auto); review=proofread-sonnet; flags=none
\begin{question}{14}
부동산 가격공시에 관한 법령상 표준지공시지가에 관한 설명으로 옳지 않은 것은?
\begin{choices}
  \item 국토교통부장관은 개별공시지가의 산정을 위하여 필요하다고 인정하는 경우에는 토지가격비준표를 작성하여 시장·군수 또는 구청장에게 제공하여야 한다.
  \item 표준지공시지가의 공시에는 표준지의 1제곱미터당 가격이 포함되어야 한다.
  \item 표준지공시지가의 공시에는 표준지에 대한 도로상황이 포함되어야 한다.
  \item 표준지공시지가에 대한 이의신청은 서면(전자문서 포함)으로 하여야 한다.
  \item 「산림조합법」에 따른 산림조합은 국유지 취득을 위해 표준지공시지가를 조사·평가할 수 있다.
\end{choices}
\end{question}
\answer{5}
